% SPDX-License-Identifier: CC-BY-4.0
% Copyright (c) 2026 Martin Schröder <info@swedishembedded.com>
\section{Limitations}\label{sec:limitations}

RSI5 did not show continuous improvement: two of four registered steps hold. It examines one task family (predicting the output
of Lua~5.1 programs) and one protocol (direct answers, thinking off); the thinking-on measurement keeps only the items every
examined model finished. It has no matched-data control, so it does not separate the evolving champion from
the growing volume of verified data; the companion paper's matched-data control could not separate them either. Four lineages
give wide intervals for steps of one or two points. The learner's gate advanced on non-significant development leads and
therefore cannot serve as a stopping rule. Every result here is for a 4B model with rank-16 adapters retrained from the base at
each generation.

The Vim family (\cref{sec:vim}) has one generation. Its arms differ in more than the registered property: G's tasks are
harder than B's and G saw each record about twice where B saw each once, so G against B does not isolate outcome-first
authoring; C against B does not separate one solution per task from fewer records. Most training records are the 27B
teacher's, and the student-only arm that would separate the teacher's contribution is not yet measured. The absolute level is
low (3.5\% of exam attempts for the best arm, against the teacher's 63\% first-attempt rate). Retention, other tools or domains,
and stability under another authoring seed of the material are not measured. The exams are 100 hand-written tasks of one
course; the per-seed intervals in \cref{fig:vim-seeds} treat attempts as independent.
